%!TEX root = ../report.tex
\documentclass[report.tex]{subfiles}
\begin{document}

\chapter{State of the Art}
\label{ch:sota}

\section{Aerial Video Segmentation Datasets}
\label{sec:sota-data}

\textbf{UAVid}~\cite{uavid2020} is an urban-drone dataset of oblique
\SI{4}{K} video, labelled densely at intervals rather than continuously. Among
its stated challenges are large scale variation, moving-object recognition and
temporal consistency across frames; the last is the property flow propagation
exploits, which is why it is the primary dataset here.

\textbf{Ruralscapes}, introduced with SegProp~\cite{segprop2020}, is a
\SI{4}{K} aerial dataset of rural scenes whose dense labels come twice as
often as UAVid's, every 50 frames. That is what makes it usable for the
single-jump propagation experiments of Section~\ref{sec:backends}.

\section{Semantic Segmentation Architectures}
\label{sec:sota-arch}

The models here follow the standard encoder-decoder pattern for dense
prediction: a backbone reduces the image to a multi-scale feature pyramid, and
a decoder fuses those scales into per-pixel logits.

The decoder throughout is \textbf{UPerNet}~\cite{upernet2018}, which combines
the Pyramid Pooling Module of \textbf{PSPNet}~\cite{pspnet2017}, adaptive
average pooling at several grid resolutions, with a Feature Pyramid Network
top-down fusion path. It is backbone-agnostic, so encoders swap without
changing the head.

Four encoder families were considered. \textbf{ConvNeXt}~\cite{convnext2022}
modernises a ResNet toward transformer design choices while remaining
convolutional, and is strongest here. \textbf{Swin}~\cite{swin2021} is a
hierarchical transformer built from shifted local-window attention, linear in
image size. \textbf{SegFormer}~\cite{segformer2021} pairs a hierarchical
transformer encoder with a minimal all-MLP decoder.
\textbf{Hiera}~\cite{hiera2023} removes architectural extras from a
hierarchical vision transformer and relies on masked-autoencoder pretraining
instead; it is the image encoder of \textbf{SAM~2}~\cite{sam2_2024}, whose
public checkpoints supply an initialisation at no training cost.

\section{Optical Flow Estimation}
\label{sec:sota-flow}

Optical flow assigns to each pixel of one frame a displacement to its
corresponding location in another. Four steps of that lineage matter here.

\textbf{Variational methods.} Brox et~al.~\cite{brox2004} formulate flow as
energy minimisation over brightness constancy, gradient constancy and a
discontinuity-preserving smoothness term, solved by nested fixed-point
iterations implementing coarse-to-fine warping. This is the classical reference
for the warping principle mask propagation rests on.

\textbf{Occlusion detection by round trip.} Following a pixel forward and then
back should return it near its origin; if it does not, the correspondence is
unreliable. Sundaram, Brox and Keutzer~\cite{sundaram2010} use this
forward-backward criterion to filter point trajectories, and it is the direct
ancestor of the validity check used here.\footnote{The criterion is often
attributed to Brox et~al.~\cite{brox2004}, which introduces the variational
warping formulation; the forward-backward occlusion test appears in the later
trajectory work~\cite{sundaram2010}.}

\textbf{Learned flow.} FlowNet~\cite{dosovitskiy2015} showed that end-to-end
CNN flow estimation was viable. \textbf{FlowNet2}~\cite{flownet2_2017} made it
competitive by stacking FlowNetC, FlowNetS and FlowNetSD into a cascade
evaluated in one forward pass; it is the backend SegProp uses, which is why it
is retained here for comparison.

\textbf{Iterative refinement.} \textbf{RAFT}~\cite{raft2020} extracts per-pixel
features, builds an all-pairs 4D correlation volume, and refines a flow field
with a recurrent unit that looks up that volume.
\textbf{SEA-RAFT}~\cite{searaft2024} simplifies and accelerates RAFT while
improving accuracy, reaching \SI{3.69}{px} endpoint error on Spring, a
\SI{22.9}{\percent} error reduction over prior published results. It is the
default backend here, though its correlation volume scales quadratically with
resolution and becomes the dominant engineering constraint at \SI{4}{K}
(Section~\ref{sec:flowbackends}).

\section{Flow-Based Label Propagation and the Gap}
\label{sec:sota-prop}

\textbf{SegProp}~\cite{segprop2020} is the closest prior work. Annotators label
roughly \SI{2}{\percent} of a video's frames and SegProp propagates those
labels to the rest along FlowNet2 flow by iterative class voting; the
propagated labels then train a segmentation network that outperforms its fully
supervised counterpart. Keyframes are selected at fixed intervals, independently
of scene content.

Two gaps follow, and together they define this project. The first is
\emph{where the starting mask comes from}. Prior work starts from a frame a
human has labelled, which is possible offline but not in flight. Here the
starting mask has to be a model prediction, so it is imperfect to begin with,
it is the expensive part of the system, and something has to decide when to
pay for it. The second is \emph{when to compute a new one}. Picking keyframes
at a fixed interval ignores what the camera is doing: one interval wastes
model calls during a hover and falls behind during fast motion. A trigger that
watches the scene instead is the alternative, and the forward-backward check
costs nothing extra, because propagation has already computed both flow
directions. Prior work also measures label quality rather than time, and time
is what decides whether propagating is worth doing at all.

\end{document}
